% 每个 frame 是一页；更多完整页面见 example/layouts/。
\section{研究问题}
\subsection{背景与目标} % 可选层级；页面标题写在 frame 后的花括号中

\begin{frame}{研究对象与目标}
  \subhead{研究问题}
  \begin{itemize}
    \item 填写研究对象和需要解决的问题。
    \item 填写已有方法的适用条件。
    \item 填写本次汇报的目标。
  \end{itemize}
  \insight{填写本页最重要的判断。}
\end{frame}

\section{研究方法}
\begin{frame}{方法与示意图片}
  \hitcolumnratio{.45} % 左栏占比；例如 .6 表示左侧较宽
  \begin{columns}[T,onlytextwidth] % T 顶部对齐，c 居中对齐
    \begin{column}{\hitleftwidth}
      \subhead{方法说明}
      填写图片所说明的机制、步骤或结果。
      \notebox{填写理解图片所需的条件。}
    \end{column}
    \begin{column}{\hitrightwidth}
      % 换成自己的图片，例如 \fig{result.png}；默认等比缩放。
      \fig{example/images/linear.png}
      \captiontext{示例图：$y=x$，$0\leq x\leq1$。来源：模板解析函数图。}
    \end{column}
  \end{columns}
\end{frame}

\section{结果与讨论}
\begin{frame}{两种方案的比较}
  \hitcolumnratio{.5}
  \begin{columns}[T,onlytextwidth]
    \begin{column}{\hitleftwidth}
      \subhead{方案一}
      \begin{itemize}
        \item 填写方法和适用条件。
        \item 填写已验证的结果。
      \end{itemize}
    \end{column}
    \begin{column}{\hitrightwidth}
      \subhead{方案二}
      \begin{itemize}
        \item 使用相同的比较维度。
        \item 说明差异及其依据。
      \end{itemize}
    \end{column}
  \end{columns}
  \insight{填写有证据支持的比较结论。}
\end{frame}
